% Pro L. Cornelio Balbo (pro-balbo) --- English translation
% Translated by Claude (Anthropic) from the Latin (Perseus).
% Style guide: STYLE.md.

\ciceroSection{1} If the authority of patrons counts in trials, the case of L. Cornelius has been defended by men of the highest standing; if the experience, by men of the greatest skill; if the talents, by the most eloquent; if the zeal, by his closest friends, joined to L. Cornelius both by his services and by the greatest familiarity. What then is my part? Of authority such as you have wished me to have, of moderate experience, of talent in no way matching my will. For to the rest by whom he has been defended, I see he owes most; how much he owes me, I shall set out in another place. At the start of my speech I lay this down: that to all who have been friends to my safety and standing, if I cannot do enough by repaying the favour, by proclaiming and acknowledging it I will at least do enough.

\ciceroSection{2} What yesterday was Cn. Pompey's gravity in speaking, gentlemen of the jury, what facility, what fullness, was being shown not by the silent thought of your minds but by visible admiration. For I have never heard anything which seemed to me more subtly said about right; nothing more wide-ranging in memory of precedents; nothing more skilled about treaties; nothing of more illustrious authority about wars; nothing weightier about the commonwealth; nothing more modest about himself; nothing more elegant about the case and the charge:

\ciceroSection{3} so that to me at last that seems to be true which some men, given to letters and the studies of learning, were thought to say as a kind of incredible thing: that to him who had thoroughly comprehended in his soul all the virtues, everything he did proceeds rightly. For what greater fullness, variety, fluency could there have been in L. Crassus, a man born for some unique faculty of speaking, had he pleaded this case, than there was in him who could give to this study only as much time as he himself, from boyhood up to this age, has rested from continuous wars and victories?

\ciceroSection{4} For which the more difficult is this concluding place of my speech. For I am succeeding a speech which has not glided past your ears, but settled deep within the minds of all, so that you may take more pleasure from the recollection of that speech than from mine, or from anyone's. But I must oblige not only Cornelius --- to whose will in his perils I cannot in any way fail --- but also Cn. Pompey, who wished me to be both the proclaimer and the advocate of his act, his judgement, his kindness, just as before you, gentlemen, recently in another case I was.

\ciceroSection{5} And to me indeed this seems worthy of the commonwealth, owed to the outstanding glory of this excellent man, proper to your own duty, sufficient to the case: that what Cn. Pompey is established to have done, all should grant him to have been at liberty to do. For nothing is truer than what he himself said yesterday --- that L. Cornelius is fighting for all his fortunes in such a way that he is called into the charge of no single offence. For he is not said to have stolen citizenship, to have falsified his birth, to have hidden in some shameless lie, to have crept onto the census-roll: one thing only is brought against him --- that he was born at Gades, which no one denies. The rest the prosecutor admits: that this man in Spain was in the most fierce war with Q. Metellus, with C. Memmius, both in the fleet and in the army; that, when Pompey came to Spain and began to have Memmius as quaestor, he never left Memmius's side; that he was besieged at Carthage; took part in those most fierce and great battles, of the Sucro and the Turia; was with Pompey to the very end of the war.

\ciceroSection{6} These things are Cornelius's own: piety toward our commonwealth, labour, perseverance, fighting, courage worthy of the highest commander, and the hope of rewards for his perils. Yet the rewards themselves are not in the act of him who received them, but of him who gave. He was therefore granted citizenship by Cn. Pompey for these reasons. The prosecutor does not deny this, but reproaches it: so that in Cornelius's case, the case is approved, but a penalty is sought; in Pompey's, the case is wounded, the penalty is none save of reputation. Thus the fortunes of a most innocent man, the act of an outstanding commander, they wish to be condemned. So Cornelius's life, Pompey's act, is brought to trial. For you grant that this man was born of the most honourable rank in that city in which he was born; and that from his earliest youth, leaving all his own affairs, in our wars he has been engaged with our commanders; that he has had no part in any toil, any siege, any battle. All these things, full of praise, are Cornelius's own; and in them there is no charge.

\ciceroSection{7} Where, then, is the charge? --- That Pompey granted him citizenship. His charge? Hardly, unless honour is to be reckoned disgrace. Whose then? In truth no one's; in the prosecutor's procedure, only of the man who granted it. Who, even if, led by favour, he had given the reward to a man less fit; even if, indeed, to a good man but not so deserving; even if, in short, anything were said to have been done not against what was lawful but against what was fitting --- yet every reproach of that sort would have to be rejected by you, gentlemen.

\ciceroSection{8} Now what is said? What does the prosecutor allege? That Pompey did what it was not lawful for him to do; which is graver than if he should say it was done not as it ought to have been. For there is something which ought not to be done, even if it is lawful; whatever is not lawful, surely ought not. Should I now hesitate, gentlemen, to plead this way: that it is not right to doubt that, what Cn. Pompey is established to have done, we should confess to have been not only lawful but actually fitting?

\ciceroSection{9} For what is wanting in this man --- which, were it present, we should think rightly attributed and granted to him? Experience? --- whose latest boyhood was the beginning of wars and the greatest commands; whose age-mates have for the most part less often seen camp than he has triumphed; who has so many triumphs as there are coasts and parts of the world, so many victories of war as there are kinds of war in nature. Or talent? --- nay, the very accidents and outcomes of things have been not the leaders but the companions of his counsels: in this one man, fortune so contended with virtue that, by the judgement of all, more was attributed to the man than to the goddess. Or modesty, integrity, religion in him --- when has diligence ever been wanting? Whom did our provinces, what free peoples, what kings, what farthest nations, see more chaste, more moderate, more holy --- nay, even hope or wish for in their own dreams?

\ciceroSection{10} What shall I say of his authority? Which is as great as in such great virtues and praises it ought to be. To this man the Senate and the Roman people gave the rewards of the highest standing without his asking, and even commands when he refused them. Of his act, gentlemen, that the question be raised in this way --- whether it was lawful for him to do what he did, or whether (I will not say it was not lawful, but) it was wickedness (for he is said to have acted contrary to the treaty, that is, contrary to the religion and faith of the Roman people) --- is this not shameful to the commonwealth, is it not to you?

\ciceroSection{11} I have heard this from my father as a boy, when Q. Metellus the son of Lucius was pleading his case on extortion --- that man, that man whose country's safety was sweeter than its sight, who chose rather to depart from the city than from his opinion --- so when he was pleading his cause, when his own ledgers were carried about for the inspection of the figures, there was not one of those gravest of men, the Roman knights serving as jurors, who did not turn his eyes away and avert himself entirely, lest perhaps it seem that anyone was doubting whether what that man had entered in his account books was true or false: shall we re-examine the decree of Cn. Pompey, gentlemen, given out under the advice of his council; shall we compare it with the laws, with the treaties, weigh everything with the bitterest diligence?

\ciceroSection{12} At Athens, they say, when a certain man who had lived holily and gravely had given testimony before them on a public matter, and, as is the Greek custom, came toward the altars to take the oath, all the jurors with one voice cried out that he should not swear. Then those Greeks, men distinguished, did not wish faith to be tied down by religion rather than by truth: shall we doubt what Cn. Pompey was, even in observing the religion of laws and treaties?

\ciceroSection{13} Do you wish him to have acted against the treaty in ignorance, or knowingly? If knowingly --- O the name of our empire! O the outstanding dignity of the Roman people! O Cn. Pompey's praise so widely and far diffused, that the home of his glory is bounded by the limits of the common empire! O nations, cities, peoples, kings, tetrarchs, tyrants --- witnesses of Cn. Pompey's not only courage in war but religion in peace! you, finally, mute regions, I implore, and the soils of the farthest lands; you, seas, harbours, islands, shores! For what is the coast, the seat, the place in which there are not the imprinted footsteps of his courage and humanity, of his spirit and his counsel? Will any man dare to say that he, gifted with an incredible and unheard-of gravity, courage, constancy, has knowingly neglected, violated, broken treaties?

\ciceroSection{14} The prosecutor obliges me with a gesture: he indicates that Pompey acted unknowingly --- as if it were lighter, when one is engaged in such great public affairs and presides over the greatest matters, to do something which one knows is not lawful, than not to know at all what is lawful! For did the man who had waged in Spain a most fierce and great war not know by what right Gades stood, or, when he knew the right of that people, did he not grasp the interpretation of the treaty? Will anyone dare to say that Cn. Pompey was ignorant of what mediocre men, men of no military experience, no soldierly study, what the very copyists profess to know?

\ciceroSection{15} I, on the contrary, hold, gentlemen, that, since in every kind and variety of arts --- even of those that are not easily learned without the highest leisure --- Cn. Pompey excels, his knowledge is unique and outstanding in the matter of treaties, of pacts, of the conditions of peoples, kings, foreign nations, and finally in the whole law of war and peace; unless perhaps those things which books teach us in shade and leisure could not be taught Cn. Pompey, when he rested by letters, nor in the doing of business by the regions themselves. And, as I judge, gentlemen, the case has been pleaded. I shall say more now from the faults of the times rather than from the kind of trial: for it is a kind of stain and slip of the age to envy virtue, to wish to break the very flower of dignity. For if Pompey had lived five hundred years ago,

\ciceroSection{16} that man, from whom the Senate as a young man and a Roman knight had often sought aid for the common safety, whose deeds had ranged through all peoples with the most distinguished victory by land and sea, whose three triumphs were witnesses that the whole circuit of the earth was held by our empire, whom the Roman people had honoured with signal and unique distinctions --- if now among us what he had done were said to have been done contrary to the treaty, who would listen? No one indeed; for since death would have extinguished envy, his deeds would have rested on the glory of an everlasting name. Will the man whose virtue, when heard of, would give no place for doubt --- when seen and looked through --- be hurt by the voice of detractors?

\ciceroSection{17} I shall therefore set Pompey aside in the rest of my speech; but you, gentlemen, hold him in your minds and memory. About the law, the treaty, the precedents, the perpetual custom of our state, I shall renew what has been said. For nothing new, nothing untouched, has been left for me to say either by M. Crassus, who has set out the whole case before you most diligently both for skill and for faith; or by Cn. Pompey, whose speech abounded with every ornament. But since, despite my refusal, both wished this last labour, as it were of finishing-off, to be applied by me to the work, I ask of you to think that I have undertaken this labour and duty out of zeal of duty rather than of speaking.

\ciceroSection{18} And before I approach the right and the case of Cornelius, something on the common condition of all of us seems to be briefly recalled, for the deflection of malice. If, gentlemen, in whatever rank of ours each was born, or in whatever fortune one was placed at the start of his being born, he must hold this state of life till old age; and if all whom either fortune has lifted or their own labour and industry has made distinguished are to be visited with punishment --- not heavier upon L. Cornelius than upon many good and brave men would the law of life and the condition seem to be set. But if many men's virtue, talent, and humanity from the lowest rank of birth and step of fortune has gained not only friendships and the means of household but the highest praise, honours, glory, dignity --- I do not see why envy should rather violate L. Cornelius's virtue than your fairness aid his modesty.

\ciceroSection{19} Therefore what is most to be asked of you I do not ask, gentlemen --- lest I seem to doubt your wisdom and humanity. But it must be asked: that you do not hate talent, are not enemies to industry, do not think humanity is to be crushed and virtue punished. This I do ask: that, if you see the case itself firm and stable in itself, you choose that the man's own ornaments be a help to the case rather than a hindrance. The case of Cornelius arises out of that law which L. Gellius and Cn. Cornelius proposed by the will of the Senate; by which law we see it has been duly ordained that those should be Roman citizens whom Cn. Pompey has, on the advice of his council, individually granted citizenship. That L. Cornelius was so granted, Pompey present says; the public registers indicate. The prosecutor admits this, but denies that anyone from a federate people could come into this citizenship unless that people had been made fundus.

\ciceroSection{20} O outstanding interpreter of right, sponsor of antiquity, corrector and emender of our state --- who attaches this penalty to treaties: that he makes federates excluded from all our rewards and benefits! For what could be said more inexperiencedly than that federate peoples must be made fundi? For that is no more proper to federates than to all free peoples. But this whole matter, gentlemen, has always been founded on this rationale and view: that, when the Roman people had ordered something, if their allies, or the Latins, accepted it, and if the same law which we held had settled itself, as in a fundus, in some other people, then by the same law that people would be held --- not so that anything would be diminished from our right, but so that those peoples would either use that right which had been established by us, or some advantage or benefit.

\ciceroSection{21} Among our forefathers C. Furius proposed a law on wills, Q. Voconius on women's inheritances; innumerable other laws on civil right have been passed; which the Latins, as they wished, accepted. The Julian law itself, by which citizenship was given to the allies and the Latins, granted no citizenship to those peoples who had not been made \textit{fundi}. On which there was great contention among the Heracleans and Neapolitans, since a great part in those communities preferred the freedom of their treaty to citizenship. In the end, this is the force both of that right and of the word: that peoples be made \textit{fundi} by our kindness, not by their own right.

\ciceroSection{22} When the Roman people has ordered something, if it is of such a kind that it seems to certain peoples, whether federate or free, must be permitted that they should themselves decide --- not concerning our affairs but concerning their own --- what right they wish to use, then it seems we must ask whether they are made \textit{fundi} or not; but concerning our commonwealth, our empire, our wars, our victory, our safety, they did not want peoples to become \textit{fundi}. And yet if it shall not be lawful for our commanders, for the Senate, for the Roman people, by setting up rewards, to draw out from the cities of allies and friends each bravest and best man to undertaking dangers for our safety, we shall have to do without the greatest help and frequently the greatest defence in dangerous and rough times.

\ciceroSection{23} But, by the immortal gods! what kind of alliance, friendship, treaty is this, that either our state should lack the Massiliot defender in its perils, lack the Gaditan, lack the Saguntine; or, if any from these peoples should arise to help our generals with toil and provisions at his own peril, who has often fought hand to hand with our enemy in the line, who has often thrown himself in the way of the enemy's weapons, of capital combat, of death --- on no terms should he be able to be rewarded with this state's prizes?

\ciceroSection{24} For on the Roman people it weighs heavily not to be able to use as allies men endowed with outstanding courage who wish to share their dangers with ours; on the allies themselves, and on those federates of whom we are speaking, it is unjust and insulting that those most loyal and most closely connected allies should be excluded from rewards and honours which lie open to stipendiaries, lie open to enemies, lie often open to slaves. For we see many stipendiaries from Africa, Sicily, Sardinia, the rest of the provinces granted citizenship; and we know that men who fled from being enemies to our generals and were of great use to our commonwealth have been granted citizenship; slaves, finally, whose right, fortune, condition is the lowest, who have well deserved of the commonwealth, we see very often publicly granted freedom --- that is, citizenship.

\ciceroSection{25} Do you, then, patron of treaties and federates, set up this condition for the Gaditans, your fellow citizens: that what is allowed to those whom we have subdued in arms with their great help, and reduced to our sway --- if the Roman people will permit, that they may be granted citizenship by the Senate, even by our generals --- this is not allowed to the Gaditans themselves? Who if they had ordained by their own decrees and laws that no one of their citizens should enter the camps of the Roman generals, that no one should expose himself to capital danger or to the hazard of life on behalf of our empire --- so that we, when we wished, might use the help of the Gaditans, but no individual man of distinguished spirit and courage should fight under our command at his own peril --- we should bear this gravely, and rightly: that the helps of the Roman people are diminished, the spirits of the bravest men weakened, that we are deprived by the eagerness of foreign men and their ancestral courage.

\ciceroSection{26} And it is no different, gentlemen, whether the federates establish these terms --- that none from those communities be allowed to come to the dangers of our wars --- or that what we have given to their citizens for their courage cannot be ratified. For we should no more use those helpers, the rewards of courage taken away, than if it were altogether forbidden to them to be in our wars. For when, on behalf of his own country, few since the human race has been have been found who without prizes set up have thrown their life upon the weapons of enemies, do you think anyone will be found who will set himself against danger on behalf of an alien commonwealth, not only with no reward set up, but even with reward forbidden?

\ciceroSection{27} But while it was very inexperienced what was said about peoples being \textit{fundi} (which is common to free peoples, not proper to federates --- from which it must be understood either that no one out of allies can become a citizen, or that he can do so even out of federates), then this our master of changing of citizenship is ignorant of every right which is set, gentlemen, not only in public laws but even in the will of private persons. For by our right neither can anyone change his citizenship unwillingly, nor, if he wishes, can he not change it --- provided he be received by that state of which he wishes to be a citizen. So if the Gaditans by name accept some Roman citizen as a Gaditan citizen, our citizen has full power to change his citizenship; nor is he hindered by the treaty from being a Gaditan citizen instead of a Roman.

\ciceroSection{28} Of two states none of us can be a citizen by civil law: he who has dedicated himself to another state can not be of this. Nor only by dedication --- as in the case of the misfortune we have seen befall the most distinguished men, Q. Maximus, C. Laenas, Q. Philippus at Nuceria; C. Cato at Tarraco; Q. Caepio, P. Rutilius at Smyrna --- whereby they became citizens of those communities, since they could not lose this one before they had turned this single soil by the change of citizenship; but also by postliminium can the change of citizenship be made. For not without reason, on Cn. Publicius Menander --- a freedman whom our forefathers' envoys, setting out for Greece, wished to take with them as interpreter --- a measure was put before the people: that this Publicius, if he should return home and from there come back to Rome, would not on that account be the less a citizen. Many in earlier memory also of their own will, uncondemned and untouched, having left these things, betook themselves into other states.

\ciceroSection{29} If then it is lawful for a Roman citizen to be a Gaditan, whether by exile, or postliminium, or by rejection of this citizenship --- to come now to the treaty, which is nothing to the case (for we are deciding the right of citizenship, not of treaties) --- what is the reason a Gaditan citizen should not be allowed to come into this state? I myself feel far otherwise. For while from all states there is a way into ours, and to our citizens lies open the way to the rest of the states, yet the more any state is joined to us by alliance, friendship, sponsion, pact, treaty, the more it seems to me to be bound by the common sharing of benefits and rewards of citizenship. And all the other states would not hesitate to receive our men into their citizenship, if we had the same right as the rest; but we cannot be of this state and of any other besides; the rest are allowed it.

\ciceroSection{30} Thus in the Greek states we see Athenians, Rhodians, Lacedaemonians, the rest from everywhere being inscribed, and the same men being of many states. By that error led, I myself have seen some inexperienced men, our citizens, at Athens, in the number of jurors and Areopagites, of a fixed tribe, of a fixed number, when they did not know that, if they had won that citizenship, they had lost this --- unless they had recovered it by postliminium. But none ever experienced in our custom and right, who wished to retain this citizenship, dedicated himself to another. But this whole topic of my disputation and speech, gentlemen, applies to the common right of changing citizenship: it has nothing peculiar to religion and treaties. For I defend the universal point: that there is no people from any region of the earth, neither so divided from the Roman people by some hatred and rupture, nor so joined in faith and goodwill, from which we are forbidden to receive or to grant citizenship.

\ciceroSection{31} O distinguished and divinely founded laws, from the very beginning of the Roman name by our forefathers established: that no one of us can be of more than one state (for the difference of states must hold a variety of right); that no one is unwillingly changed in citizenship, nor remains in citizenship unwillingly! For these are the firmest foundations of our liberty: that each is master both of holding and of giving up his own right. But this without doubt most of all founded our empire and increased the name of the Roman people: that that founder of this city, Romulus, by his treaty with the Sabines, taught that this state should be increased even by receiving enemies; on whose authority and example the largesse and sharing of citizenship was never broken off by our forefathers. So both from Latium many --- as Tusculans, Lanuvians --- and from the rest of kinds whole peoples were received into citizenship, as of the Sabines, Volscians, Hernicians; from which states neither would they have been forced to change citizenship if they had not wished to, nor, if any had attained our citizenship by the kindness of the Roman people, would their treaty have seemed violated.

\ciceroSection{32} For some treaties exist --- as of the Cenomani, Insubres, Helvetii, Iapydes, and some likewise of the Gallic barbarians --- in whose treaties it has been excepted that no one of them should be received as our citizen. But if exception makes it not lawful, then where it is not excepted, there it must be lawful. Where then is it in the Gaditan treaty that the Roman people should receive no Gaditan as a citizen? Nowhere. And if anywhere it were, the Gellian and Cornelian law, which had given Pompey definite power of granting citizenship, would have removed it. ``The treaty,'' he says, ``has been excepted, \textit{if anything is sacrosanct}.'' I forgive you, if you neither know the laws of the Carthaginians (for you had left your own state) nor could you inspect our laws --- for they themselves drove you back from their cognizance by a public trial.

